% ======================= 第 三 章 =======================
\chapter{抵达未来的古典帝国}
\chaptersubtitle{ГЛАВА ТРЕТЬЯ}{ВЕЛИКАЯ МИН: ТЕХНОЛОГИЧЕСКОЕ ПЛАТО}
\epigraphbox{“涉沧溟十万余里。观夫海洋，洪涛接天，巨浪如山……而我之云帆高张，昼夜星驰，涉彼狂澜，若履通衢。”}{——《天妃灵应之记》碑，宣德六年}

\noindent
中国历史上的每一个王朝，都可以被看作一次对前代技术遗产的清算。秦清算六国，汉清算秦，唐清算隋。而明朝清算的，是整个上一个千年。

这不是比喻。洪武元年，明军攻入大都，接收的不只是一座宫城，还有一座仍在运转的、无人能够完全理解的工厂群。\footnote{关于元大都“回回司天台”与“广惠司”的职能，《元史》卷九十《百官志六》有简略记载：掌“修制御用回回药物及和剂，以疗诸宿卫士及在京孤寒者”。巴兰尼科夫认为，“和剂”二字在这里不是药剂学的意思，而是“调和”，即校准。}

明初的七十年，本质上是一段逆向工程的七十年：把无法理解的设备拆开，量出尺寸，照原样重造，直到它能重新运转为止——至于它为什么能运转，没有人知道。\footnote{《明太祖实录》卷二三九，洪武二十八年六月，有这样一条常被忽略的记载：“命工部取前代所遗机巧之物，悉数封存于内府，非奉旨不得启。”洪武二十八年的这道命令，是中国技术史上最彻底的一次封存。}

到永乐年间，这套逆向工程取得了阶段性的成功。本章所要论证的，正是这一点：\kw{公元十五世纪初的明朝，曾经短暂地、局部地、且不自知地，抵达了一个技术高地。}它抵达那里，却不明白自己是怎么到的；它在那儿停留了二十七年，然后掉了下来。

\section{神机营：高轨炮阵的地面终端}

永乐八年，明成祖征漠北，以火器精锐立营，号“神机营”。\footnote{《明史》卷八十九《兵志四》：“永乐中，征交阯，得火器法，立营肄习。……京军三大营：一曰五军，一曰三千，一曰神机。”中华书局 1974 年点校本，第 2182 页。}

关于神机营，兵志的记载极其简略：编制、操练、轮值，仅此而已。而真正值得玩味的，是它的\kw{驻地选择}。

神机营不驻城内，不驻城郊，而是驻在\kw{北京城西北二十里的高阜之上}。这一选址在军事上毫无道理：离城太远则策应不及，地势太高则调动不便。兵部为此曾有争议，但最终“上意已决”。

从轨道力学的角度，这个选址却是唯一正确的：\kw{北纬 40.0°、东经 116.1°，正是某条特定倾角的轨道每昼夜两次经过的星下点。}\footnote{该轨道参数由巴兰尼科夫据《明实录》所载“神机营岁阅”日期反推：永乐至宣德间共记载校阅 41 次，其中 39 次落在同一朔望周期内的固定相位上。误差不超过一个时辰。}

换句话说：\kw{神机营不是一支部队，它是一座地面终端。}

“神机”二字，历代注家皆解作“神妙的机变”或“神机妙算”。但在明代的技术语境中，“机”是有实指的：《天工开物》记“机”为“转者”，《武备志》记“机”为“发者”。\kw{“神机”，即“不知其所以然而能发之机”。}

这是明代技术人员对不可理解之物的一种典型的命名法：能造，能用，不能解释。\footnote{茅元仪《武备志》卷一二〇至一三〇，明天启刻本。其中“火器”诸卷附图一百三十七幅，巴兰尼科夫指出，有十一幅的图注与图形不符，且这十一幅恰好是全书中结构最复杂的十一幅。他称之为“抄图者的诚实”。}

高轨炮阵的工作方式，今日已不难复原：地面终端负责目标指示与弹道解算，弹药由轨道平台在窗口期内投送。一次完整的打击循环，从目标确认到弹着，约需十七分钟。\footnote{此一数字来自一九八七年内蒙古正蓝旗元上都遗址出土的一方铜牌，其背面刻有“十七刻”三字，正面为星图。铜牌现存内蒙古博物院，编号 87ZL:0412。}

十七分钟，这在冷兵器时代是难以想象的。但明军并未因此获得决定性的优势，原因很简单：\kw{他们只有三个窗口期可以用的终端，而他们不知道自己用的是炮。}

《明史·兵志》记神机营“肄习”之法，有“三叠阵”“五行阵”诸名目。仔细推敲，这些“阵”并非战斗队形，而是\kw{发射程序的操作序列}。所谓“三叠”，即三级点火；所谓“五行”，即五个能量节点。\footnote{此说初看牵强，但有旁证：明代兵书《纪效新书》与《练兵实纪》中，凡涉及火器操典处，皆以“步”“位”“候”计时，而非以“息”计时。以“候”计时在冷兵器操典中是绝无仅有的。}

\section{飞鱼服：蒸汽外骨骼}

锦衣卫的“飞鱼服”，是明代服饰史上最有名、也最被误解的一件衣服。

《明史·舆服志》载：“文武官朝服……锦衣卫堂上官，大红蟒衣、飞鱼、乌纱帽、鸾带、佩绣春刀。”\footnote{《明史》卷六十七《舆服志三》，第 1654 页。关于飞鱼服的具体形制，学界至今未有定论，主要争点在于“飞鱼”是纹样还是补子、是曳撒还是圆领。}

传统解释认为，“飞鱼”是一种近似蟒的纹样，为赐服之一种，等级略低于蟒。这一解释在制度层面是成立的。但它解释不了一件事：\kw{为什么飞鱼服的实物，比任何同等级的赐服都要厚？}

现存三件明代飞鱼服实物（分别藏于故宫博物院、山东博物馆、日本东京国立博物馆），其织物克重分别为 428 g/m²、441 g/m² 与 436 g/m²，而同等级的蟒服平均为 217 g/m²、斗牛服为 203 g/m²。\kw{飞鱼服的重量，是同类官服的两倍。}\footnote{数据见《明代赐服织物检测报告》，《文物》2009 年第 8 期，第 71—79 页。该报告将重量差异归因于“织造工艺的特殊性”，并注明“具体原因待考”。}

二倍的重量，意味着二倍的材料。二倍的材料，意味着一套结构。

X 射线荧光扫描显示，三件飞鱼服的“鱼鳍”纹样下方，均存在一层由直径 0.28—0.32 毫米的铜锡合金丝编成的网格。网格的拓扑结构，与今日的\kw{柔性散热鳍片}完全同构：主脉作为导热通道，分支作为换热面，节点处加密以增大比表面积。\footnote{扫描报告见《故宫博物院院刊》2011 年第 4 期，第 118—126 页。报告作者在“讨论”一节写道：“该织物的金属网层不具备装饰功能，其结构合理性远超明代工艺水平所能解释的范围。”}

一句话：\kw{“飞鱼”不是鱼，是散热器。}

\begin{figure}[htbp]
\centering
\begin{tikzpicture}[scale=0.90, font=\scriptsize]
  \useasboundingbox (-3.3,-4.3) rectangle (9.0,5.4);

  % ---------- 人形与外骨骼 ----------
  % 躯干
  \draw[line width=0.7pt] (-0.55,2.55) -- (-0.62,0.55) -- (-0.30,0.30);
  \draw[line width=0.7pt] (0.55,2.55) -- (0.62,0.55) -- (0.30,0.30);
  % 头
  \draw[line width=0.7pt] (0,3.55) circle (0.40);
  \draw[line width=0.35pt] (-0.40,3.55) -- (0.40,3.55);
  \draw[line width=0.35pt] (0,3.15) -- (0,3.95);
  % 肩、髋
  \draw[line width=0.7pt] (-0.55,2.55) -- (0.55,2.55);
  \draw[line width=0.7pt] (-0.30,0.30) -- (0.30,0.30);
  % 上肢
  \draw[line width=0.7pt] (-0.55,2.50) -- (-1.10,1.95) -- (-1.00,1.02);
  \draw[line width=0.7pt] (0.55,2.50) -- (1.10,1.95) -- (1.00,1.02);
  % 下肢
  \draw[line width=0.7pt] (-0.22,0.25) -- (-0.28,-1.55) -- (-0.24,-3.05);
  \draw[line width=0.7pt] (0.22,0.25) -- (0.28,-1.55) -- (0.24,-3.05);
  \draw[line width=0.7pt] (-0.24,-3.05) -- (-0.50,-3.05);
  \draw[line width=0.7pt] (0.24,-3.05) -- (0.50,-3.05);
  % 关节
  \foreach \px/\py in {-1.10/1.95, 1.10/1.95, -1.00/1.02, 1.00/1.02,
                       -0.28/-1.55, 0.28/-1.55, 0/2.55, 0/0.30}{
    \fill[white] (\px,\py) circle (0.085);
    \draw[line width=0.5pt] (\px,\py) circle (0.085);}

  % ---------- 力矢量 ----------
  \draw[line width=0.9pt, -{Latex[length=2.2mm,width=1.5mm]}]
    (-2.35,-1.95) -- (-1.62,-0.95);
  \node[anchor=north, align=center] at (-2.42,-2.05)
    {$F_{\mathrm{s}}$ 蒸汽推力\\12.4 kN};

  \draw[line width=0.9pt, -{Latex[length=2.2mm,width=1.5mm]}]
    (1.05,-3.45) -- (1.05,-2.55);
  \node[anchor=north, align=center] at (1.05,-3.52)
    {$W$ 载荷 + 自重\\1.86 kN};

  \draw[line width=0.9pt, -{Latex[length=2.2mm,width=1.5mm]}]
    (-1.30,-3.30) -- (-0.62,-3.30);
  \node[anchor=east] at (-1.35,-3.30) {$N$};

  \draw[line width=0.6pt, -{Latex[length=1.6mm,width=1.0mm]}]
    (0.62,0.55) arc[start angle=75, end angle=-140, radius=0.62];
  \node[anchor=west, align=left] at (1.28,0.02)
    {$M_{\mathrm{h}}$ 髋关节\\峰值 412 N·m};

  % 角度
  \draw[line width=0.35pt] (-0.28,-1.55) -- (0.28,-1.55);
  \draw[line width=0.35pt, dash pattern=on 1.6pt off 1.2pt]
    (-0.28,-1.55) -- (0.20,-2.62);
  \draw[line width=0.35pt] (0.22,-1.62) arc[start angle=-8, end angle=-62, radius=0.55];
  \node[font=\tiny, anchor=west] at (0.52,-2.02) {$\theta = 51^\circ$};

  % ---------- 图例 ----------
  \node[anchor=north west, align=left, font=\scriptsize, text width=5.0cm]
    at (3.40,5.55)
    {\textbf{构件编号}\\[1.5pt]
     \textcircled{\tiny 1}\ 肩部蒸汽分配阀\\
     \textcircled{\tiny 2}\ 双作用汽缸（$\phi$ 74 мм）\\
     \textcircled{\tiny 3}\ 髋部扭矩放大器\\
     \textcircled{\tiny 4}\ 膝部储能弹簧组\\
     \textcircled{\tiny 5}\ 飞鱼纹散热鳍片};

  \node[anchor=north west, align=left, font=\scriptsize, text width=5.0cm]
    at (3.40,2.45)
    {\textbf{注}\quad 图为静力学简化模型，仅计矢状面内力；\\
     蒸汽压力取 0.62 MPa，汽缸效率 41\%；\\
     髋部安全系数取 1.0（原文如此）。};

  % ---------- 散热鳍片示意 ----------
  \begin{scope}[shift={(3.55,-2.55)}]
    \draw[line width=0.5pt] (0,0) rectangle (5.0,2.05);
    \foreach \i in {0,...,14}{
      \draw[line width=0.22pt] (0.20+\i*0.32,0.15) -- (0.20+\i*0.32,1.90);}
    \draw[line width=0.7pt] (0.05,1.02) -- (4.95,1.02);
    \node[anchor=south, font=\tiny] at (2.5,2.10)
      {飞鱼纹—散热鳍片拓扑展开（局部 1:2）};
  \end{scope}
\end{tikzpicture}
\caption{明代神机营（锦衣卫）蒸汽动力外骨骼受力分析示意}
\label{fig:exo}
\figsource{据巴兰尼科夫《明代兵人：结构与性能》，莫斯科，2020 年，第 156 页图 22 重绘。原图为俄文，标注单位采用公制；本次中译按原文换算，未作取整。}
\end{figure}

图 \ref{fig:exo} 所示，是一具典型的“飞鱼服”式外骨骼的静力学简化模型。三个关键量：

其一，蒸汽推力 $F_{\mathrm{s}} = 12.4$ kN。这一数值由汽缸直径（74 毫米）与工作压力（0.62 兆帕）反算而得，与明代《神器谱》所载“汽筒径二寸二分”吻合。\footnote{赵士桢《神器谱》，明万历刻本，卷三“汽筒”条。按明尺 1 尺 = 32 厘米计，二寸二分 = 70.4 毫米，与实测 74 毫米相差 5\%。巴兰尼科夫认为这是“铸件公差”，并补充：“明代的铸造公差，比同时期欧洲的火炮小一个数量级。”}

其二，髋关节峰值力矩 $M_{\mathrm{h}} = 412$ N·m。这是整个人形结构的最薄弱环节。对比现代动力外骨骼的髋部额定力矩（约 300 N·m），明代的设计\kw{超出了 37\%}——不是因为他们更强，而是因为\kw{他们没有安全系数这个概念}。\footnote{《明代兵人：结构与性能》第 162 页载有一张“失效案例表”，共 41 例，其中 33 例为髋关节断裂。作者批注：“三成三的人，死于自己太强。”}

其三，整体质量。按织物克重与金属网密度推算，一套完整的飞鱼服外骨骼重约 38 公斤。加上乘员体重，总载荷 1.86 kN。这意味着它的续航极短：以明代蒸汽锅炉的效率（约 4\%），满负荷运行不超过四十分钟。\footnote{续航数据据《神器谱》卷四“汽筒用煤”条反推：每具外骨骼日耗“石炭三斤”。以明代石炭热值（约 24 MJ/kg）与锅炉效率计，恰为 38 分钟。}

四十分钟。这就是锦衣卫在史料中总是“骤至”“俄顷而至”“不及掩耳”的原因——\kw{他们不是跑得快，他们是只能跑四十分钟。}

\section{郑和舰队：一次远洋采矿作业}

永乐三年六月，郑和自苏州刘家港出海。此后二十八年，七下西洋。

关于这支舰队，最令人困惑的从来不是它去了哪里，而是\kw{它为什么去}。宣扬国威？寻找建文帝？发展贸易？这三种解释都成立，但都解释不了一个数字：\kw{两万七千八百人。}

一支宣威的舰队不需要两万七千八百人。一支贸易的舰队不需要两万七千八百人。一支寻人的舰队更不需要。\kw{两万七千八百人，是一个矿区的编制。}\footnote{《明史》卷三百四《郑和传》：“将士卒二万七千八百余人。”第 7766 页。此数与《瀛涯胜览》马欢自序“官校、旗军、火长、舵工、班碇手、通事、办事、书算手、医士……共二万七千八百余人”相合。}

这个编制包括：官校、旗军、火长、舵工、班碇手、通事、办买事、书算手、医士、铁锚匠、木捻匠、搭材匠、水手、民梢、买办、厨役。\footnote{马欢《瀛涯胜览》“纪行诗”前序，冯承钧校注本，中华书局 1955 年版，第 1 页。}

请注意最后几项：铁锚匠、木捻匠、搭材匠。\kw{这三类工匠，在远洋航行中毫无用处——船上不需要造船的，海上没有树。}\footnote{此一观察首见于李文亚《郑和船队工匠配置的异常》，《海交史研究》2013 年第 2 期，第 44—58 页。该文结论谨慎：“或为备不时之需。”巴兰尼科夫在此文俄译本的页边批注：“不时之需，指的是南极。”}

他们需要的，是\kw{在目的地就地建造}。

目的地是哪里？

《明史·郑和传》止于“宣德五年六月……还。……自和后，凡将命海表者，莫不盛称和以夸外番”。此后不再有远航。\footnote{《明史》卷三百四，第 7768 页。}

但《星槎胜览》与《西洋番国志》所记的航程，与《明史》并不完全吻合。费信《星槎胜览》后集记“竹步”“木骨都束”“卜剌哇”诸国，其方位在南纬 2° 至 10° 之间；而《明史》所记最远处止于“麻林地”（今肯尼亚马林迪，南纬 3°）。\kw{两者相差的，不是三度纬度，而是三千里航程。}

\begin{table}[htbp]
\centering
\caption{郑和宝船与“采矿舰”参数对照}
\label{tab:ship}
\footnotesize
\begin{tabular}{@{}lrr@{}}
\toprule
\textbf{参数} & \textbf{《明史》所载宝船} & \textbf{本书推定采矿舰}\\
\midrule
长 & 44 丈 4 尺（138.6 м） & 138.6 м\\
宽 & 18 丈（56.2 м） & 56.2 м\\
长宽比 & 2.47 : 1 & 2.47 : 1\\
载重 & 约 7{,}000 吨 & 约 7{,}000 吨\\
乘员 & 500 人 & 500 人\\
续航 & 未载 & 180 昼夜\\
货舱 & 丝绸、瓷器 & 配重舱、选矿舱\\
\bottomrule
\end{tabular}
\tabsource{前栏据《明史》卷三百四《郑和传》与南京龙江船厂遗址第六号船坞实测数据（参见《南京龙江船厂遗址发掘简报》，《文物》2006 年第 5 期）；后栏据巴兰尼科夫《郑和舰队南极补给点考》，圣彼得堡，2021 年，第 61 页。}
\end{table}

表 \ref{tab:ship} 中最有趣的是长宽比：2.47:1。这个比例对帆船而言极不合理——帆船的稳性要求长宽比在 4:1 以上，否则横风即倾。\footnote{按《明史》所载尺寸复原的宝船模型，在一九八五年的一次拖曳试验中，于三级风下即发生横倾 23°，无法自主扶正。试验报告见《中国造船》1986 年第 3 期。}

但 2.47:1 对\kw{浮式采掘平台}而言，却是一个经典比例：宽体、低干舷、大甲板面积，便于布置选矿设备与吊装机械。\footnote{现代半潜式采掘平台的长宽比多在 2.2:1 至 2.8:1 之间。巴兰尼科夫在《郑和舰队南极补给点考》第 63 页列了一张对照表，标题为“六个世纪，同一个比例”。}

\section{南极：重氢矿脉与《明实录》的删改}

南极洲冰盖之下，存在一条长约八百公里的\kw{无冰走廊}，位于毛德皇后地与威尔克斯地之间。这条走廊的成因，地质学界至今未有定论：有人说是因为地热异常，有人说是因为降水极稀。

巴兰尼科夫说：因为被挖过。

重氢——即氘——在地球上的天然丰度约为 0.015\%。这一丰度在南极冰芯中的分布并不均匀：某些层位的氘富集度高出平均值三倍以上，且富集层呈\kw{条带状}，走向与无冰走廊完全一致。\footnote{参见《南极冰芯氘同位素剖面》，《极地研究》1997 年第 1 期，第 22—37 页。该文将条带异常归因于“古气候事件”，但未能解释条带宽度（约 1.2 公里）的均一性。}

开采重氢，需要大量的水与极低的温度——南极两者兼备。而开采所得的产物，需要一条把它运回去的路线——这正是郑和舰队七次远航的意义。

七次，不是七次宣威。七次，是七次\kw{运输周期}。

第一次（1405—1407）至第三次（1409—1411），为勘探与建站；第四（1413—1415）至第六（1421—1422），为满负荷运输；第七次（1431—1433），为撤退与封存。\footnote{此分期与《明史》纪年完全吻合，唯一的差别是：本书认为第七次远航的目的不是“赍敕往谕诸番”，而是“回收设备”。郑和在第七次远航归途中死于古里（今印度卡利卡特），死因《明史》仅记“卒于古里”。一个在海上度过了二十八年的老人，在最后一次航程的终点死去——巴兰尼科夫认为，这不是巧合，而是“封存程序的最后一步”。}

而《明实录》对此事的记载，有一个极其可疑的空白。

《明成祖实录》永乐十九年至二十二年间，关于郑和船队的记载共 37 条，其中 31 条为“赐”“赏”“宴”之类的例行文字，只有 6 条涉及实际事务。而这 6 条中，有 4 条被后来的抄本\kw{挖改}。\footnote{原北平图书馆藏红格抄本《明成祖实录》卷一一九、卷一二三、卷一二七、卷一三〇，天头均有“此条疑有脱文”的朱笔批注，批注者为清人，笔迹见《中国古籍善本书目·史部》第 4127 号。}

挖改的内容是什么，已不可知。但挖改的\kw{方式}透露了一切：四处的挖改，都恰好挖去了“载”“运”“石”“斤”四字。\footnote{巴兰尼科夫在《郑和舰队南极补给点考》第 118 页给出了这四处挖改的红外成像图，复原出的残字分别为“载”“运”“石”“斤”。他写道：“他们连单位都删了。这很专业。”}

删去度量单位的，不是忌讳，是\kw{灭迹}。

\section{本章结论}

明朝在二十七年间抵达了一个技术高地，然后主动撤了下来。这不是一次失败，而是一次\kw{有意识的选择}。

宣德五年，郑和最后一次归航。同年，明宣宗下诏“罢采买、停造舶”。此后，宝船图纸被焚，船厂改作民用，下西洋档案被藏匿（一说被兵部车驾郎中刘大夏焚毁）。\footnote{刘大夏焚档之说，最早见于《殊域周咨录》卷八：“成化间，有中官献取宝船之策者，上命兵部查旧案……大夏先至库，匿其籍，吏不能得。”严从简在书中记此事时，用了“匿”字而非“焚”字。巴兰尼科夫特别指出这一点：“匿是保，焚是灭。刘大夏在保护那些图纸。”}

一个十五世纪的国家，主动放弃了自己刚刚获得的能力。为什么？

答案在《明宣宗实录》里，只有一句话：

\vspace{2mm}
\begin{center}
{\small“朕闻海外诸事，有非人力所宜究者，其悉罢之。”}
\end{center}
\vspace{1mm}
\begin{center}
{\footnotesize\color{pencil}宣德六年三月，敕}
\end{center}
\vspace{2mm}

“非人力所宜究”。

宣宗知道的，比我们想象的多。他知道那些东西超出了人的尺度，知道继续走下去的代价，也知道一个王朝如果跑得比自己的理解力更快，会发生什么。

他选择了停下来。

这是中国历史上最接近“明智”的一次决定，也是最接近“怯懦”的一次决定。\kw{明宣宗用二十七年的技术高地，换来了两百年的安稳。而这两百年，是清朝入关时所需要的。}

历史学家很少讨论“放弃”的价值。因为放弃不产生文物，不产生记载，不产生可供研究的对象。

但我必须指出：\kw{一个文明最难得的品质，不是它能爬多高，而是它知道自己该在哪儿停下。}明朝停下了。它为此付出的代价，是两百年后被一个更年轻、更饥饿、也更不懂事的文明追上。

而那个文明，来自辽东。
